\documentclass[11pt,a4paper]{article}
\usepackage[margin=25mm]{geometry}
\usepackage{amsmath,amssymb,amsthm,hyperref}
\setlength{\emergencystretch}{2em}
\newtheorem{lemma}{Lemma}
\newtheorem{theorem}[lemma]{Theorem}
\newcommand{\gen}[1]{\langle#1\rangle}
\title{A perfect split from two support bounds}
\author{Human Compression project: a component of the L019 recovery}
\date{}
\begin{document}\maketitle
Let $G=\gen{\rho_0,\ldots,\rho_{r-1}}$ be a transitive
string action by involutions on $n$ points. Assume that its
noncommutation diagram is connected, every maximal omission
is intransitive, and $2r>n+2$.
An eligible edge of colour $i$ joins different orbits of the
actual maximal omission $G_i$. If there is exactly one such
edge, deleting it from the full permutation graph gives the
two $G_i$-orbits. Call this cut perfect if, after choosing
names $A,B$ for these orbits, $G_{<i}$ fixes $B$ pointwise
and $G_{>i}$ fixes $A$ pointwise. The generator $\rho_i$
may still have internal edges in both sides.

We use three complete earlier arguments: paired orbit cuts,
including the bound $2r\le n$ when every colour has two
eligible edges; the orbit-budget and head-support theorems
in their stated parity-free form; and the ambient interval
inequality. The last says that an intransitive subgroup
generated by a nonempty contiguous interval $I$, each of
whose colours has two ambient eligible edges, has support
$U$ satisfying $2|I|\le |U|-c$, where $c\ge1$ is its
number of nontrivial orbits. No eligibility relative to a
smaller subgroup is substituted for this hypothesis.

\begin{lemma}[Two separated unique cuts]\label{three}
Suppose unique cuts at $i<j$ have edges in series. Write their
three regions as $A,M,D$, so the $i$-omission orbits are
$A$ and $B=M\cup D$, while the $j$-omission orbits are
$C=A\cup M$ and $D$. Suppose $j\ge i+2$, $G_{<i}$ fixes
$B$, and $G_{>j}$ fixes $C$. If $h>i$ acts on $A$ and
$g<j$ acts on $D$, then
\[
                         n\ge 2r+h-g.
\]
The assertion also allows $i=0$ or $j=r-1$.
\end{lemma}
\begin{proof}
Apply the outside-side bound at $i$ to get
$|A|\ge i+h+1$. Apply it with all indices reversed at $j$:
the crossing index is $r-1-j$ and the active index is
$r-1-g$. Thus $|D|\ge2r-j-g-1$.

Let $b\in M$ be the endpoint of the $i$-edge towards $j$,
and $c\in M$ the endpoint of the $j$-edge towards $i$.
A simple path from $b$ to $c$ in $M$ avoids colour $i$'s
unique crossing and stays in $B$. The no-gap argument at
$i$ applies to any such path to a point moved by $\rho_j$:
it must contain every index $i+1,\ldots,j-1$.
Equivalently, if some such index $k$ were missing, commute
the letters of index greater than $k$ to the beginning.
They all fix $b$, while $\rho_j$ fixes $b$ and commutes
with the remaining lower-index word. It would fix $c$,
which it does not. Hence the path has at least $j-i-1$
edges and $|M|\ge j-i$. Summing the three disjoint region
bounds proves the assertion.
\end{proof}

\begin{theorem}\label{perfect}
Under the standing hypotheses there is a perfect split.
\end{theorem}
\begin{proof}
The paired-cut bound ensures a unique eligible edge for
at least one colour. Suppose, for a contradiction, that
no unique cut is perfect. The orbit-budget theorem then
provides a shared mark on one side of some cut; reverse
the entire index order if necessary to put it to the right.
In this fixed order choose a largest cut index $i$ having
some shared mark of index greater than $i$.
The primitive interval side furnished by the orbit-budget
argument must lie on the right, since every shared mark
acts on that side. Thus name the sides so that $G_{<i}$
fixes $B$ and $G_{>i}$ is transitive on $B$. Both sides
have size at least two: an interior singleton contradicts
the singleton budget, while an endpoint singleton is already
a perfect split.

Let $h>i$ be the largest shared index. The head theorem
gives $h<r-1$ and, with $X=\operatorname{supp}(G_{>h})$,
\begin{equation}\label{x}
 X\subseteq B\setminus\{b\},\qquad 2h\le n-|X|-1.
\end{equation}
At $i=0$ we choose the names of the sides according to the
final tail, as expressly allowed by that theorem; the
condition on $G_{<i}$ remains automatic.
If every colour greater than $h$ had two ambient eligible
edges, the interval theorem, applied to the tail fixing $A$,
would give $2(r-1-h)\le |X|$. Together with \eqref{x}
this would imply $2r\le n+1$, a contradiction.

Choose the least $j>h$ having a unique eligible edge.
Since $G_{>h}$ is supported in $B$, this edge lies in $B$.
The two unique edges are graph bridges, so deleting both
partitions the connected full graph into precisely three
regions $A,M,D$, in that order. Thus the $j$-orbits are
$C=A\cup M$ and $D$. The region $M$ contains the two
endpoints facing one another, and $j\ge i+2$.
By maximality of $i$, no mark greater than $j$ is shared
at this new cut. Since this cut is not perfect, it has a
shared mark less than $j$. Let $g$ be its first shared index.
Again $|C|,|D|\ge2$.

We claim that $G_{>j}$ fixes $C$, and that the reversed
head theorem has its final tail in $C$. If $j<r-1$, the
primitive interval at $j$ lies to the left, since there
is a shared mark to the left and none to the right.
If this primitive side were $D$, then $G_{>j}$ would
fix $D$. Apply the last-mark bound with indices reversed,
with first side $C$ and primitive side $D$. Its crossing
index $r-1-j$ is positive. The original mark $\rho_0$
acts nontrivially on $C$: for $i>0$ it has its nonempty
support in $A\subset C$, and for $i=0$ it contains the
$i$-bridge, which lies wholly in $C$. The last-mark bound
would give $n\ge2r$, impossible. Therefore the primitive
side is $C$, proving that $G_{>j}$ fixes it. The reversed
head theorem now places its final tail in $C$.

If $j=r-1$, $G_{>j}$ is empty and fixes $C$ automatically.
The reversed last-mark bound excludes $g=0$, since that
would make the original mark $\rho_0$ shared at this
endpoint cut. The final reversed tail acts on a single
side, and it contains $\rho_0$, which acts nontrivially
on $C$ as just shown. Its side is therefore $C$.
In both cases, with $Y=\operatorname{supp}(G_{<g})$,
the reversed head bound is exactly
\begin{equation}\label{y}
 Y\subseteq C,\qquad 2(r-1-g)\le n-|Y|-1.
\end{equation}
The reversed index is $r-1-g$, not $r-g$.

Lemma~\ref{three} applies, because $h$ acts on $A$ and
$g$ on $D$. It gives $n\ge2r+h-g$. Thus $g\le h$
is impossible, and $g>h$.
The two groups defining $X,Y$ together contain every
mark. Since $G$ is transitive on more than one point,
no point is fixed by all its generators. Consequently
\begin{equation}\label{union}
                         X\cup Y=\Omega.
\end{equation}
Every colour strictly between $h$ and $g$ has two ambient
eligible edges, by the minimal choice of $j$ and $g<j$.
Its subgroup has support $V\subseteq X\cap Y$ and fixes
both $i$-omission sides setwise, so is intransitive.
If this interval is nonempty, the ambient interval bound
gives $2(g-h-1)\le |V|\le|X\cap Y|$. If it is empty,
the same weak inequality is immediate. Summing it with
\eqref{x} and \eqref{y}, and then using \eqref{union},
gives
\[
\begin{split}
 2r&=2h+2(g-h-1)+2(r-1-g)+4\\
   &\le(n-|X|-1)+|X\cap Y|+(n-|Y|-1)+4\\
   &=n+2,
\end{split}
\]
contradicting the standing strict inequality. This proves
the theorem.
\end{proof}

\paragraph{Relation to L019 and attribution.}
Cameron--Fernandes--Leemans, \emph{The number of string
C-groups of high rank}, arXiv:2212.12723v2, establish
perfect-split existence in Section~5. The fracture strategy,
unique crossings and support overlap remain attributed to
that source and the earlier alternating-group work.
The present assembly uses the complete acquired ambient
interval proof and the newly reworked actual-orbit support
bounds. Its separated-cut argument retains all weak equality
cases; the strict inequality $2r>n+2$ is what excludes them.
No parity-specific graph classification or quotient IP is
silently imported. The global middle-colour exchange in the
acquired paired-cut theory retains its separately disclosed
root contribution and distinct review.

For a high-rank symmetric C-marking, the accepted initial
fracture theorem supplies the standing omission hypothesis
in its stated degree range; its small-degree qualification
must still be retained when assembling the principal source.
This theorem alone does not establish a pure end mark,
the degree/rank bijection, or the six enumerative constants.
No whole-source credit, new conjecture, Lean declaration or
compression measurement is claimed.
\end{document}
